\documentclass[10pt,a4paper]{amsart}
\usepackage[margin=19mm]{geometry}
\usepackage{amsmath,amssymb,mathtools}
\usepackage[scr=rsfs,cal=euler]{mathalfa}
\usepackage{xfrac}
\usepackage[unicode,hidelinks,bookmarks=false]{hyperref}
\newcommand{\ie}{i.e.\ }
\newcommand{\cf}{cf.\ }
\newcommand{\resp}{respectively\ }
\newcommand{\Subsection}{\subsection}
\providecommand{\nobreakdash}{\nobreak\hbox{-}}
\setlength{\emergencystretch}{2em}
\multlinegap0pt
\usepackage{bm}
\usepackage{bbm}
\usepackage{dsfont}
\usepackage{xspace}
\usepackage{cancel}
\usepackage{mathtools}
\usepackage{amsthm}
\makeatletter
\@ifundefined{theorem}{\newtheorem{theorem}{Theorem}}{}
\makeatother
\title[Further results on the lower bound on reduced Zagreb index o]{Further results on the lower bound on reduced Zagreb index of trees}
\author{Milan Ba\v{s}i\'c, Aleksandar Ili\'c}
\date{Source: arXiv:2604.06044v1 (CC BY 4.0)}
\begin{document}
\maketitle

\noindent\emph{Source and licence.} Further results on the lower bound on reduced Zagreb index of trees. Version arXiv:2604.06044v1 (CC BY 4.0), distributed under the Creative Commons Attribution 4.0 International licence, as recorded in the arXiv licence metadata for this exact version and shown on the abstract page https://arxiv.org/abs/2604.06044v1. Full rights notice: licenses/arxiv-2604.06044-CC-BY-4.0.txt.\par\smallskip

\noindent\textit{Editorial scope.} Counted unit: the Theorem whose source label is \texttt{None} in the article (source0.tex lines 680--683, source label \texttt{None}), delivered together with the article's own complete original proof of it (source0.tex lines 684--749, one \texttt{proof} environment of 66 lines). No mathematics is written, repaired or re-derived: statement and proof are the article's own text line for line. Required context retained uncounted: none. Every definition, display and label of those lines is retained unchanged. Omitted from this package and not redistributed: 1--679, 750--967. Nothing inside a retained excerpt is omitted or altered: every retained body is delivered exactly as the article's lines are written, so no omission receipt of any kind accompanies this package. Citations: the retained excerpts carry no citation command at all, so no bibliography entry of the article is transcribed and no reference is redistributed. Reference adaptation: none was needed and none was made; every reference inside the delivered unit resolves inside it, so no unresolved reference or citation is produced. Printed slips kept literal: no printed word, symbol, hypothesis or number of the counted statement or of its proof was altered. Layout and class: the article's own class is \texttt{elsarticle} with packages including natbib, graphicx, multirow, amsmath, amssymb, amsfonts, amsthm, mathrsfs, appendix, xcolor, textcomp, manyfoot, booktabs, algorithm; the wrapper is the lane's pinned \texttt{amsart} editorial wrapper at 10pt on A4, which restates the theorem-like environments the retained text needs and adds the article's own macro definitions that the retained text needs (none), with extra wrapper packages (bm, bbm, dsfont, xspace, cancel, mathtools). The delivered numbering is the unit's own. Assets: no retained span contains a figure environment, a graphics-inclusion command, a PDF-page inclusion, a table or a TikZ picture; no image file is copied into this package, the package declares no asset dependency and no image file is redistributed with it. Licence: version arXiv:2604.06044v1 of this article is distributed under the Creative Commons Attribution 4.0 International licence, as recorded in the arXiv licence metadata for this exact version and shown on the abstract page https://arxiv.org/abs/2604.06044v1. A full rights notice is delivered at licenses/arxiv-2604.06044-CC-BY-4.0.txt. Characters: every retained excerpt is pure ASCII, so the delivered engine, XeLaTeX with its default Latin Modern text font, reports no missing character.
\begin{theorem}
Let $T$ denote a tree with maximum degree $3$ and order $n\geq 7$. 
It follows that $GRM_{-2}(T)\geq -\lfloor \frac{n-1}{3} \rfloor-2=-(k+2)$. %The equality holds if $T\equiv T^1_{\text{opt}}(k)$ or $T\equiv T^2_{\text{opt}}(k)$ or $T\equiv T^3_{\text{opt}}(k)$.
\end{theorem}

\begin{proof}
The following equations hold for any tree $T$ with a maximal degree of $3$

\begin{eqnarray}
    &&n_1 + n_2 + n_3  = n \nonumber\\ 
    &&n_1 + 2n_2 + 3n_3 = 2n-2 \nonumber\\ 
    &&m_{12} + m_{13}  = n_1 \label{1st}\\
    &&m_{12} + 2m_{22} + m_{23}  = 2n_2 \nonumber \\ 
    &&m_{13} + m_{23} + 2m_{33}  = 3n_3 \label{second}\\  \nonumber
\end{eqnarray}



Using the previous system, we observe that we can express the variables $n_1$, $n_2$, $m_{12}$, $m_{13}$, and $m_{33}$ in terms of $n_3$, $m_{22}$, and $m_{23}$.

\begin{eqnarray}
    n_1 &=& 2 + n_3\label{2nd}\\
    n_2 &=& n - 2 - 2n_3\label{second'}\\
    m_{33}&=&n-n_3-3-m_{22} - m_{23} \label{third}\\
    m_{13}&=&5n_3+2m_{22} + m_{23}-2n+6 \label{4th}\\
    m_{12}&=&2n-4n_3-2m_{22} - m_{23}-4\label{5th}.\\ \nonumber
\end{eqnarray}

By subtracting   (\ref{third}) from (\ref{4th}) we can deduce that 
\begin{eqnarray}
    m_{13}-m_{33}&=& 6n_3-3n+3m_{22}+2m_{23}+9 \label{6th}.  \\ \nonumber
 \end{eqnarray}

From (\ref{5th}),  it can be concluded that 
$m_{23}\leq  2n-4n_3-2m_{22} -4$, under the assumption that $m_{12} \geq 0$.
Combining this inequality with equation (\ref{6th}) we deduce that
\begin{eqnarray}
m_{13}-m_{33}\leq n-2n_3-m_{22}+1.\label{7th}
\end{eqnarray}
Now let $n=3k+r$, for $1\leq r\leq 3$.

\medskip

Suppose that $n_3\geq \lfloor \frac{n-1}{3} \rfloor+1=k+1$.

Using (\ref{7th}) and the fact that $n_3\geq k+1$ we obtain that 
\begin{eqnarray}
\label{inequality2}
m_{13}-m_{33}&\leq& (3k+r)-2(k+1)-m_{22}+1\leq k+r-1\leq k+2.\\\nonumber
\end{eqnarray}
The equality holds if and only if $m_{12}=m_{22}=0$ and $n_3=k+1$.

\medskip

Suppose that $n_3\leq \lfloor \frac{n-1}{3} \rfloor=k$.

According to (\ref{1st}) and (\ref{2nd}) we get that
\begin{eqnarray}
\label{inequality1}
    m_{13}-m_{33}&\leq & n_1=n_3+2\leq k+2.\\\nonumber
\end{eqnarray}

%Given that $m_{13} - m_{33} = -GRM_{-2}(T)$, it can be established that $T_{\text{opt}}^1(k)$, $T^2_{opt}(k)$, and $T^3_{opt}(k)$ are trees, with their parameters $m_{13}$ and $m_{33}$ satisfying the aforementioned inequality. This is %considered in light of the respective relations (\ref{Zagreb_opt}), (\ref{Zagreb_opt_2}), and (\ref{Zagreb_opt_3}).

%Therefore these trees are optimal in the class of all trees of the order $n$ and maximal degree equal to $3$, as indicated by the inequality
%\begin{eqnarray*}
%GRM_{-2}(T)&=&-(m_{13} - m_{33})\geq -(k+r-1)\geq -(k+2)\\
%&=&GRM_{-2}(T_{\text{opt}}^1(k))=GRM_{-2}(T^2_{opt}(k))=GRM_{-2}(T^3_{opt}(k)),
%\end{eqnarray*}
%for any tree $T$ such that $n_3\geq \lfloor \frac{n-1}{3} \rfloor+1$.
\end{proof}

\end{document}
